\section{Conclusão e trabalhos futuros}

Reconstruímos do ChEMBL\_37 uma base de 27.714 moléculas, cada uma associada a um de quatro alvos (EGFR, BACE1, DRD2 e HDAC1), e treinamos uma árvore de decisão para prever o alvo. Com divisão agrupada, os cinco descritores físico-químicos dão acurácia de \accGrupDesc\%; acrescentar 12 contagens de caracteres do SMILES leva a \accGrupAmbos~$\pm$~\sdGrupAmbos\% (ganho relativo de \ganhoAmbos\%), e as contagens sozinhas já chegam a \accGrupSmiles\%. A divisão aleatória superestima o resultado em \gapMin\% a \gapMax\% (relativo) porque cerca de um quarto das moléculas de teste tem gêmea de descritores no treino. TPSA e o número de átomos aromáticos são os atributos mais importantes, e o par HDAC1--EGFR concentra 28,7\% dos erros.

Na comparação com a Random Tree e com conjuntos de árvores, a árvore de decisão segue como a melhor árvore única. A Random Tree fica \relRT\% abaixo dela quando ajustada e \relRTweka\% abaixo com os padrões do Weka, e varia mais com a semente: o sorteio de atributos esconde o melhor corte de cada nó, e uma árvore só não tem como compensar. Nos conjuntos, o mesmo sorteio ajuda: torna as árvores diferentes entre si, e 300 Random Trees juntas, a Random Forest, acertam 6,7 pontos a mais que a árvore de decisão na seed 42. O \emph{gradient boosting} ajustado é o modelo mais preciso (\accGBaj\%, \relGBaj\% acima da árvore). A Random Forest chega a \accRF\% sem ajuste e é a única cuja confiança separa acertos de erros. Em troca, só a árvore de decisão oferece regras legíveis.

Na base de \rnLinhas{} pedidos de empréstimo, uma rede neural com padronização e parada antecipada acerta \rnAccRede~$\pm$~\rnSdRede\%: supera a regressão logística, um neurônio só, nas 10 seeds, e a regra de um corte em \rnSeedsRedeRegra{} delas. Sem padronização, ela não aprende (\rnSemPadAcc\%). Com a paciência padrão de 10 épocas e taxa de aprendizado de $10^{-3}$, pararia cedo demais (\rnParadaAccLentaDez\% de validação). Mesmo bem ajustada, a rede perde para a árvore de decisão (\rnAccArv\%) em todas as seeds, treina cerca de \rnRazaoTreino{} vezes mais devagar e exige uma busca de hiperparâmetros \rnRazaoBusca{} vezes mais longa. O motivo é a forma da regra que gerou os rótulos, um degrau no score de crédito: a árvore o representa com uma pergunta, e a rede só o aproxima com uma rampa, onde ficam seus erros. Nesta base, a rede funciona, mas não compensa.

Quatro caminhos seguem do que foi visto. Representações estruturais da molécula (\emph{fingerprints}), calculadas com uma biblioteca de química, enxergariam subestruturas como o grupo C(=O)NO, que as contagens não veem. Uma divisão por esqueleto molecular daria uma estimativa menos otimista para moléculas novas. A curva de aprendizado ainda sobe, o que sugere buscar mais dados no ChEMBL, com outros alvos ou um limiar de pChEMBL mais baixo. Por fim, calibrar as probabilidades do \emph{gradient boosting}, hoje confiante demais, permitiria usar a confiança do modelo mais preciso para escolher quais previsões revisar. Na base de empréstimos, o próximo passo seria repetir a comparação em dados de crédito reais, onde a relação com o score é gradual, e separar a variação causada pelos pesos iniciais da rede da variação causada pela divisão dos dados.
